Experiments were run on a MacBook Pro with an Apple M4 Pro processor and 24 GB of RAM, running macOS 26.6.1. The NVIDIA CUDA toolkit is not supported on macOS. Training was on CPU. The software packages and their respective versions used for the data pipeline, model training and evaluation are listed in Table~\ref{tab:software}.

\begin{table}[htbp]
\centering
\begin{tabular}{ll}
\toprule
Package & Version \\
\midrule
Python & 3.12.13 \\
PyTorch & 2.13.0 \\
NumPy & 2.5.2 \\
pandas & 3.0.5 \\
polars & 1.43.2 \\
pyarrow & 25.0.1 \\
DuckDB & 1.5.5 \\
SciPy & 1.18.0 \\
Matplotlib & 3.11.1 \\
\bottomrule
\end{tabular}
\caption{Software packages and versions}
\label{tab:software}
\end{table}

The TTF futures data from Bloomberg and the options data from Databento are licensed and cannot be distributed. The processing steps are documented in Sec.~\ref{sec:data-pipeline} for futures data and in Sec.~\ref{sec:pricing-setup} for options data. The pipeline is reproducible, given access to the same sources.

%\todo{A fixed seed was set for the Python, NumPy and PyTorch random number generators at the start of each run. Bit-level reproducibility is nevertheless not guaranteed on the MPS backend, as not all operations used are deterministic. Repeated runs under an identical seed differed by [X]. All distributional statistics are therefore reported as means across seeds, following Sec. 3.5. A single training run of [N] epochs took approximately [X] minutes, and the experiments reported in Ch. 4 comprise [N] runs in total.}

%\todo{Implementation. The implementation is based on the reference Fin-GAN code published by \textcite{vuletic2024fingan}. The model architecture and training loop were adopted from that implementation and adapted for the MPS backend and for the univariate TTF return series. The data pipeline described in Sec. 3.1, the evaluation procedures in Sec. 3.5 and the option pricing setup in Sec. 3.6 were implemented independently. do I NEED TO SUBMIT THE CODE?}